\documentclass[preview]{standalone}

\usepackage{amsmath}
\usepackage{amssymb}
\usepackage{bettelini}
\usepackage{stellar}
\usepackage{definitions}

\begin{document}

\id{measuretheory-borel-sets}
\genpage

\section{Borel sets}

\begin{snippetdefinition}{borel-sigma-algebra-metric-space-definition}{Borel \(\sigma\)-algebra}
    Let \((X,\tau)\) be a \topologicalspace. The \emph{Borel \(\sigma\)-algebra}
    on \(X\) is
    \[
        \mathcal B(X)=\sigma(\tau),
    \]
    the \(\sigma\)-algebra generated by the open subsets of \(X\). Its elements
    are called \emph{Borel sets}.
\end{snippetdefinition}

\begin{snippet}{borel-sigma-algebra-basic-sets}
    The \borelsigmaalgebra contains every open set and every closed set, as well
    as many sets that are neither open nor closed.
\end{snippet}

\begin{snippetdefinition}{g-delta-set-definition}{\(G_\delta\) set}
    Let \(X\) be a \topologicalspace. A subset of \(X\) is a
    \emph{\(G_\delta\) set} if it is a countable intersection of open sets.
\end{snippetdefinition}

\begin{snippetdefinition}{f-sigma-set-definition}{\(F_\sigma\) set}
    Let \(X\) be a \topologicalspace. A subset of \(X\) is an
    \emph{\(F_\sigma\) set} if it is a countable union of closed sets.
\end{snippetdefinition}

\begin{snippetproposition}{g-delta-f-sigma-sets-are-borel}{\(G_\delta\) and \(F_\sigma\) sets are Borel}
    Every \(G_\delta\) set and every \(F_\sigma\) set in a \topologicalspace
    \(X\) belongs to \(\mathcal B(X)\).
\end{snippetproposition}

\begin{snippetexample}{real-intervals-g-delta-f-sigma-example}{Intervals as \(G_\delta\) and \(F_\sigma\) sets}
    Every interval in \(\realnumbers\), with any choice of open or closed
    endpoints, is both a \(G_\delta\) set and an \(F_\sigma\) set. For example,
    if \(a<b\), then
    \[
        [a,b)
        =\bigcap_{n=1}^{\infty}\left(a-\frac1n,b\right)
        =\bigcup_{n=1}^{\infty}
        \left[a,b-\frac{b-a}{n+1}\right].
    \]
\end{snippetexample}

\begin{snippet}{countable-subsets-of-r-are-f-sigma}
    Every countable subset of \(\realnumbers\) is an \(F_\sigma\) set because
    singletons are closed. In particular, \(\rationalnumbers\) is an
    \(F_\sigma\) set.
\end{snippet}

\begin{snippettheorem}{baire-category-theorem}{Baire category theorem}
    A nonempty complete \metricspace cannot be written as a countable union of
    closed sets with empty interior. Equivalently, if
    \[
        X=\bigcup_{n=1}^{\infty}F_n
    \]
    and every \(F_n\) is closed, then at least one \(F_n\) has nonempty interior.
\end{snippettheorem}

\begin{snippetexercise}{rationals-not-g-delta-exercise}{The rationals are not a \(G_\delta\) set}
    Prove that \(\rationalnumbers\) is not a \(G_\delta\) subset of
    \(\realnumbers\).

    \emph{Hint.} Use the Baire category theorem.
\end{snippetexercise}

\begin{snippetsolution}{rationals-not-g-delta-exercise-solution}{The rationals are not a \(G_\delta\) set}
    Suppose that
    \[
        \rationalnumbers=\bigcap_{n=1}^{\infty}G_n
    \]
    with every \(G_n\) open. Since \(\rationalnumbers\) is dense and is
    contained in each \(G_n\), every \(G_n\) is dense. Thus each
    \(\realnumbers\difference G_n\) is closed with empty interior.
    If \(\{q_n\}_{n=1}^{\infty}\) enumerates \(\rationalnumbers\), then
    \[
        \realnumbers
        =\bigcup_{n=1}^{\infty}(\realnumbers\difference G_n)
        \union\bigcup_{n=1}^{\infty}\{q_n\}.
    \]
    This writes the complete metric space \(\realnumbers\) as a countable union
    of closed sets with empty interior, contradicting the Baire category theorem.
\end{snippetsolution}

\section{Borel hierarchy}

\begin{snippet}{borel-hierarchy-introduction}
    Countable intersections of \(F_\sigma\) sets form the class
    \(F_{\sigma\delta}\), while countable unions of \(G_\delta\) sets form
    \(G_{\delta\sigma}\). Alternating these operations produces further levels
    such as \(G_{\delta\sigma\delta}\). No fixed finite number of alternations
    exhausts all Borel sets; in general, the Borel hierarchy continues through
    the countable ordinals.
\end{snippet}

\begin{snippetproposition}{generated-sigma-algebra-cardinality-continuum}{Cardinality of a generated \(\sigma\)-algebra}
    Let \(\mathcal S\subseteq\powerset(X)\) and suppose that
    \(\cardinality{\mathcal S}=\cardinality{\realnumbers}\). Then
    \[
        \cardinality{\sigma(\mathcal S)}=\cardinality{\realnumbers}.
    \]
\end{snippetproposition}

\begin{snippetproof}{generated-sigma-algebra-cardinality-continuum-proof}{generated-sigma-algebra-cardinality-continuum}{Cardinality of a generated \(\sigma\)-algebra}
    Every element of \(\sigma(\mathcal S)\) can be represented by a countable
    well-founded expression built from members of \(\mathcal S\) using
    complements and countable unions. Such expressions can be encoded by
    countable labelled trees. There are at most
    \[
        \cardinality{\realnumbers}^{\aleph_0}
        =\cardinality{\realnumbers}
    \]
    such codes, so \(\cardinality{\sigma(\mathcal S)}\leq
    \cardinality{\realnumbers}\). The reverse inequality follows from
    \(\mathcal S\subseteq\sigma(\mathcal S)\).
\end{snippetproof}

\begin{snippetproposition}{non-borel-subsets-of-r-exist}{Non-Borel subsets of \(\realnumbers\)}
    There exist subsets of \(\realnumbers\) that are not Borel measurable.
\end{snippetproposition}

\begin{snippetproof}{non-borel-subsets-of-r-exist-proof}{non-borel-subsets-of-r-exist}{Non-Borel subsets of \(\realnumbers\)}
    The open intervals generate \(\mathcal B(\realnumbers)\), and the family of
    open intervals has the cardinality of \(\realnumbers\). Hence
    \[
        \cardinality{\mathcal B(\realnumbers)}
        =\cardinality{\realnumbers}.
    \]
    Cantor's theorem gives
    \(\cardinality{\powerset(\realnumbers)}>
    \cardinality{\realnumbers}\), so not every subset of \(\realnumbers\) is Borel.
\end{snippetproof}

\begin{snippetexercise}{no-countably-infinite-sigma-algebra-exercise}{A \(\sigma\)-algebra cannot be countably infinite}
    Prove that every infinite \(\sigma\)-algebra has at least the cardinality of
    the continuum. In particular, no \(\sigma\)-algebra is countably infinite.
\end{snippetexercise}

\begin{plainhtml}
    The Borel sigma algebra is not large enough for every construction needed in Lebesgue integration. Its completion will lead to the Lebesgue sigma algebra.
\end{plainhtml}

\end{document}
